\chapter{Testowanie systemu}

Proces weryfikacji poprawności działania systemu został podzielony na kilka etapów, obejmujących testy jednostkowe warstwy backendowej, testy integracyjne \gls{API} oraz manualne testy interfejsu użytkownika. Takie podejście pozwoliło na wykrycie błędów zarówno na poziomie logicznym, jak i funkcjonalnym.

\section{Testy jednostkowe i integracyjne backendu}
Warstwa serwerowa aplikacji, zaimplementowana w języku Python, została pokryta testami automatycznymi przy użyciu frameworka \texttt{pytest}. Głównym celem tych testów była weryfikacja poprawności działania poszczególnych endpointów \gls{API} oraz logiki biznesowej generatora danych.

Testy, zlokalizowane w katalogu \texttt{backend/tests}, obejmują następujące scenariusze:
\begin{itemize}
    \item Weryfikacja autentykacji: Sprawdzenie poprawności generowania i walidacji tokenów \gls{JWT}. Testy upewniają się, że dostęp do chronionych zasobów bez ważnego tokenu jest blokowany (kod błędu 401 Unauthorized).
    \item \gls{CRUD} projektów: Testowanie operacji tworzenia, odczytu i usuwania projektów. Weryfikowana jest poprawność zapisu struktury \gls{JSON} do bazy danych PostgreSQL.
    \item Walidacja schematów Pydantic: Sprawdzenie, czy system poprawnie odrzuca błędne żądania (np. brak wymaganych pól, niepoprawne typy danych).
    \item Symulacja generowania danych: Uruchamianie silnika generującego dla prostych konfiguracji w celu potwierdzenia, że zwracane dane są zgodne z zadeklarowanym typem.
\end{itemize}

Przykładowy test weryfikujący działanie endpointu \texttt{/generate} wygląda następująco:
\begin{verbatim}
def test_generate_endpoint(client, auth_headers):
    payload = {
        "config": {"job_name": "Test Job", "rows": 10},
        "tables": [...]
    }
    response = client.post("/generate", json=payload, headers=auth_headers)
    assert response.status_code == 200
    data = response.json()
    assert "status" in data
    assert data["status"] == "success"
\end{verbatim}

\section{Testy manualne interfejsu użytkownika}
Ze względu na dynamiczny charakter interfejsu, duży nacisk położono na testy manualne. Scenariusze testowe obejmowały pełne User Journeys:

\subsection*{Scenariusz 1: Utworzenie schematu sklepu internetowego}
Celem tego testu było sprawdzenie integralności relacyjnej generowanych danych.
\begin{enumerate}
    \item Utworzenie tabeli \texttt{users} z polami \texttt{id}, \texttt{name}, \texttt{email}.
    \item Utworzenie tabeli \texttt{orders} z polem \texttt{user\_id} jako klucz obcy do tabeli \texttt{users}.
    \item Uruchomienie generowania dla 100 użytkowników i 500 zamówień.
    \item Weryfikacja: Sprawdzenie, czy każde zamówienie przypisane jest do istniejącego użytkownika. Test zakończył się sukcesem - silnik poprawnie rozwiązał zależności w 100\% przypadków.
\end{enumerate}

\subsection*{Scenariusz 2: Integracja z LLM}
Testowano zachowanie systemu przy dużym obciążeniu zapytaniami do modelu językowego.
\begin{enumerate}
    \item Zdefiniowanie pola typu \texttt{llm} z promptem generującym opis produktu.
    \item Ustawienie liczby wierszy na 50.
    \item Obserwacja: Początkowo czas generowania wyniósł średnio 1.66 sekundy na rekord przy użyciu modelu GPT-4o-mini, co przy większych zbiorach danych było wartością niesatysfakcjonującą. W celu optymalizacji wprowadzono mechanizm równoległego przetwarzania zapytań (batch processing) z wykorzystaniem \texttt{asyncio.gather}. Po tej zmianie czas generowania serii 50 rekordów spadł do 11 sekund (ok. 0.22s na rekord), co stanowi ponad 7-krotne przyspieszenie. System poprawnie obsłużył asynchroniczność, nie blokując interfejsu użytkownika podczas oczekiwania na odpowiedź z \gls{API} OpenAI.
\end{enumerate}

\section{Analiza wydajności}
Przeprowadzono testy wydajnościowe mające na celu określenie limitów systemu. Testy wykonano na laptopie deweloperskim z kartą graficzną NVIDIA RTX 5060, przy użyciu lokalnego modelu Ollama (Llama 3). Jako schemat testowy wykorzystano \texttt{ecommerce\_llm\_schema} zawierający 4 tabele, w tym tabelę recenzji z polem \texttt{review\_text} generowanym przez \gls{LLM}.

\begin{table}[H]
    \centering
    \begin{tabular}{|l|c|c|c|}
        \hline
        \textbf{Typ danych}       & \textbf{Liczba wierszy} & \textbf{Czas (Faker)} & \textbf{Czas (\gls{LLM})}   \\
        \hline
        Proste (Imię, Adres)      & 1\,000                  & 0,8s                  & --                          \\
        \hline
        Proste                    & 10\,000                 & 5,2s                  & --                          \\
        \hline
        Złożone (z kontekstem AI) & 100                     & --                    & 52,24s (po optymalizacji)   \\
        \hline
        Złożone (z kontekstem AI) & 10\,000                 & --                    & 8\,038s ($\approx$2h 14min) \\
        \hline
    \end{tabular}
    \caption{Zestawienie czasów generowania danych}
    \label{tab:performance_tests}
\end{table}

Jak widać w Tabeli \ref{tab:performance_tests}, generowanie danych przy użyciu algorytmów pseudolosowych (Faker) jest niezwykle szybkie i skalowalne liniowo. Wąskim gardłem systemu jest generowanie danych przy użyciu \gls{LLM}, co potwierdza słuszność decyzji o zastosowaniu przetwarzania wsadowego z kontrolą współbieżności. Szczegółowy opis procesu optymalizacji algorytmu generowania opisano w Rozdziale~\ref{sec:llm_optimization}.

\section{Testy End-to-End (\gls{E2E})}
W celu zapewnienia najwyższej jakości oprogramowania oraz weryfikacji integralności całego systemu, wprowadzono kompleksowe testy typu \gls{E2E}. Testy te mają na celu symulację rzeczywistych scenariuszy użytkowania aplikacji, sprawdzając współpracę wszystkich komponentów systemu - od interfejsu użytkownika w przeglądarce, przez warstwę \gls{API}, aż po logikę biznesową i bazy danych.

Do realizacji testów \gls{E2E} wybrano framework Playwright. Jest on nowoczesnym narzędziem umożliwiającym automatyzację testów w wielu przeglądarkach (Chromium, Firefox, WebKit) oraz zapewniającym wysoką niezawodność dzięki mechanizmom automatycznego oczekiwania na elementy interfejsu.

Zaimplementowany scenariusz testowy pokrywa kluczową ścieżkę krytyczną\cite{happy_path} użytkownika:
\begin{enumerate}
    \item Autoryzacja: Scenariusz rozpoczyna się od wejścia na stronę logowania. Automat testowy wprowadza dane uwierzytelniające administratora i weryfikuje poprawne przekierowanie do głównego widoku aplikacji.
    \item Definicja schematu: Wirtualny użytkownik dodaje nową tabelę o nazwie ,,users'' oraz konfiguruje jej strukturę, dodając kolumnę ,,name''. Test weryfikuje poprawność interakcji z formularzami oraz dynamiczne aktualizacje interfejsu.
    \item Generowanie danych: Następuje uruchomienie procesu generowania danych poprzez kliknięcie przycisku ,,Run Generation''. Test monitoruje asynchroniczny proces, oczekując na komunikaty o postępie przekazywane przez WebSocket.
    \item Weryfikacja wyników: Ostatnim krokiem jest potwierdzenie zakończenia zadania sukcesem (pojawienie się komunikatu ,,Data generated!'') oraz sprawdzenie dostępności przycisku pobierania wyników.
\end{enumerate}

Dzięki zastosowaniu symulowania odpowiedzi backendu w niektórych testach, możliwe jest izolowane testowanie warstwy frontendowej, co przyspiesza proces deweloperski i uniezależnia testy interfejsu od dostępności zewnętrznych usług czy modeli \gls{LLM}. Takie podejście gwarantuje, że błędy w warstwie prezentacji są wykrywane natychmiastowo, zanim kod trafi na środowisko produkcyjne.

\section{Testy bezpieczeństwa}
W celu zapewnienia wysokiego poziomu bezpieczeństwa, przeprowadzono serię testów z wykorzystaniem narzędzi klasy \gls{SCA} oraz \gls{SAST}.

\subsection{Analiza podatności zależności (\gls{SCA})}
Do identyfikacji znanych podatności w wykorzystywanych bibliotekach użyto narzędzi \texttt{npm audit} (dla warstwy frontendowej) oraz \texttt{pip-audit} (dla warstwy backendowej).

\begin{itemize}
    \item \textbf{Frontend}: Analiza początkowo wykazała obecność podatności w bibliotece \texttt{axios} (wersja $\le$ 1.3.4). W ramach działań naprawczych, zaktualizowano bibliotekę do najnowszej stabilnej wersji, eliminując zidentyfikowane zagrożenia.
    \item \textbf{Backend}: Audyt zależności Pythonowych nie wykazał krytycznych podatności w głównych bibliotekach produkcyjnych (FastAPI, Uvicorn, SQLAlchemy).
\end{itemize}

\subsection{Statyczna analiza kodu (SAST)}
Kod źródłowy części serwerowej został poddany analizie przy użyciu narzędzia \textbf{Semgrep} \cite{semgrep}.
Wykryto potencjalne zagrożenie w konfiguracji mechanizmu \gls{CORS}, polegające na zbyt szerokim uprawnieniu dostępu (\texttt{allow\_origins=["*"]}).
Problem ten został rozwiązany poprzez wprowadzenie konfiguracji opartej na zmiennych środowiskowych, co pozwala na ścisłe zdefiniowanie zaufanych domen w środowisku produkcyjnym:
\begin{verbatim}
    allow_origins=os.getenv("ALLOWED_ORIGINS").split(",")
\end{verbatim}



